\documentclass[10pt]{article}
\usepackage[T1]{fontenc}
\usepackage[utf8]{inputenc}
\usepackage[french]{babel}
\usepackage{amsmath,amssymb,mathrsfs,mathpazo}
\usepackage{geometry}
\geometry{a4paper,margin=2cm}
\usepackage{graphicx,array,tabularx,tabularray,colortbl,multirow,multicol,booktabs,enumitem,pifont,ulem,fancybox,diagbox,cancel}
\usepackage{pstricks,pstricks-add,pst-plot,pst-tree,pst-node,pst-text,pst-func,pst-math,pst-bspline,pst-3dplot,pst-eucl,pst-solides3d}
\let\pscalculate\undefined
\usepackage{pst-bezier}
\usepackage{tikz,pgfplots,tkz-tab}
\usetikzlibrary{arrows,arrows.meta,patterns,calc,positioning,trees,decorations.markings,shapes.geometric,shapes.misc,shapes.arrows,intersections,decorations.pathmorphing,decorations.pathreplacing,backgrounds,shadows,babel}
\pgfplotsset{compat=1.18}
\usepgfplotslibrary{fillbetween}
\usepackage[np]{numprint}
\usepackage{listings,makecell,pas-tableur,twoopt}
\let\barre\undefined
\usepackage[pstricks]{tablvar}
\definecolor{darkgreen}{RGB}{0,100,0}
\providecommand{\red}{\color{red}}
\providecommand{\blue}{\color{blue}}
\providecommand{\green}{\color{green}}
\providecommand{\black}{\color{black}}
\providecommand{\white}{\color{white}}
\providecommand{\gray}{\color{gray}}
\providecommand{\cyan}{\color{cyan}}
\providecommand{\magenta}{\color{magenta}}
\providecommand{\yellow}{\color{yellow}}
\providecommand{\violet}{\color{violet}}
\providecommand{\orange}{\color{orange}}
\providecommand{\euro}{\texteuro}
\providecommand{\EUR}{\texteuro}
\providecommand{\ieme}{\textsuperscript{e}}
\providecommand{\degres}{\ensuremath{^\circ}}
\providecommand{\degree}{\ensuremath{^\circ}}
\providecommand{\vv}[1]{\overrightarrow{#1}}
\providecommand{\Coord}[1]{#1}
\providecommand{\decosix}{}
\providecommand{\starredbullet}{\ensuremath{\bullet}}
\pagestyle{empty}
\setlength{\parindent}{0pt}
\usepackage{adjustbox}\DeclareUnicodeCharacter{2236}{:}\DeclareUnicodeCharacter{2190}{\ensuremath{\leftarrow}}
\providecommand{\R}{\mathbb{R}}
\providecommand{\N}{\mathbb{N}}
\providecommand{\Z}{\mathbb{Z}}
\providecommand{\Q}{\mathbb{Q}}
\providecommand{\D}{\mathbb{D}}
\providecommand{\C}{\mathbb{C}}
\providecommand{\e}{{\mathrm{e}}}
\providecommand{\vect}[1]{\overrightarrow{#1}}
\providecommand{\vectt}[1]{\overrightarrow{\mathrm{#1}}}
\providecommand{\vvt}[1]{\overrightarrow{\mathrm{#1}}}
\providecommand{\np}[2][]{{#2~#1}}
\providecommand{\vv}[1]{{\overrightarrow{#1}}}
\renewcommand{\d}{{\mathrm{d}}}
\providecommand{\degres}{{^\circ}}
\providecommand{\degre}{{^\circ}}
\providecommand{\blue}{\color{blue}}
\providecommand{\red}{\color{red}}
\providecommand{\enskip}{\quad}
\providecommand{\llbracket}{\lbrack\!\lbrack}
\providecommand{\rrbracket}{\rbrack\!\rbrack}
\providecommand{\euro}{{\text{€}}}
\providecommand{\no}{{\text{n°}}}
\providecommand{\interval}[2]{{[#1;#2]}}
\providecommand{\intervFO}[2]{{[#1;#2[}}
\providecommand{\dotfill}{\ldots}

\providecommand{\ds}{\displaystyle}
\providecommand{\textding}[1]{#1}
\providecommand{\Oij}{\(\mathrm{O};\vec{\imath},\vec{\jmath}\)}
\providecommand{\Oijk}{\(\mathrm{O};\vec{\imath},\vec{\jmath},\vec{k}\)}
\providecommand{\Ouv}{\(\mathrm{O};\vec{u},\vec{v}\)}
\providecommand{\intervFF}[2]{\ensuremath{[#1;#2]}}
\renewcommand{\euro}{\ifmmode\text{\texteuro}\else\texteuro\fi}
\DeclareUnicodeCharacter{00B0}{\ifmmode^{\circ}\else\textdegree\fi}
\DeclareUnicodeCharacter{20AC}{\ifmmode\text{\texteuro}\else\texteuro\fi}
% Énoncés : Éducation nationale. Transcription et corrigés : APMEP, Denis Vergès et contributeurs. Découpage et classement : CoopMaths, Sébastien Lozano, Éric Elter et contributeurs. Adaptation de la présentation : Maths Explorer.
% CC BY-SA : https://coopmaths.fr/www/about/
% https://coopmaths.fr/alea/static/e3c/2020/tex/e3c_2020_00_sujet20_1_cor.tex
\begin{document}
\noindent\textit{APMEP / CoopMaths. CC BY-SA. Présentation adaptée par Maths Explorer.}\par\medskip
\subsection*{Question 1}
	
	Soit \(ABC\) un triangle tel que \(AB = 6\), \(AC = 3\) et \(\widehat{ BAC} = \dfrac{\pi}{3}\).
	
	On a \(\overrightarrow{AB} \cdot \overrightarrow{AC} = AB \times AC \times \cos(\widehat{BAC}) = 6 \times 3 \times \dfrac{1}{2} = 3 \times 3 = 9\).
	
	\subsection*{Question 2}
	
	\[
	\dfrac{f(1+h) - f(1)}{h} = h^2 + 3h - 1 \quad \text{ou} \quad \dfrac{f(1+h) - f(1)}{(1+h)-1} = h^2 + 3h - 1.
	\]
	
	On a par définition
	\[
	\lim_{h \to 0} \dfrac{f(1+h) - f(1)}{(1+h)-1} = f'(1) \quad \text{si cette limite existe}.
	\]
	
	Or,
	\[
	\lim_{h \to 0} h^2 + 3h - 1 = -1.
	\]
	
	Donc \(f'(1) = -1\) (nombre dérivé de la fonction \(f\) en 1).
	
	\subsection*{Question 3}
	
	Soit \(f\) la fonction définie sur \(\mathbb{R}\) par \(f(x) = (x + 2)e^x\).
	
	\(f\) est un produit de fonctions dérivables sur \(\mathbb{R}\), donc sur \(\mathbb{R}\),
	\[
	f'(x) = 1e^x + (x + 2)e^x = e^x (1 + x + 2) = (x + 3)e^x.
	\]
	
	\subsection*{Question 4}
	
	Soit \(f\) une fonction telle que \(f(2) = 5\) et \(f'(2) = -1\). Dans un repère, la tangente à la courbe représentative de \(f\) au point d’abscisse 2 a pour équation :
	
	On sait qu’une équation de la tangente à la courbe représentative de \(f\) au point d’abscisse 2 est :
$y  = f'(2)(x - 2)+f(2)$, soit  $y - 5 = -1(x - 2)$ ou encore $y = -x + 7$.
	
	
	\subsection*{Question 5}
	
	
	Pour l’équation réduite de la tangente à la courbe \(C_f\) au point \(A\), l’ordonnée à l’origine est celle de \(B\) soit \(-\dfrac{5}{3}\) et son coefficient directeur est celui de la droite \((AB)\), soit
	\[
	\dfrac{-\dfrac{5}{3} - \dfrac{4}{3}}{0 - 1} = \dfrac{-3}{-1} = 3.
	\]
	
	On sait que ce coefficient directeur est égal au nombre dérivé \(f'(1) = 3\).

\end{document}
